%% annex/p04/universal-reductions-supplement.tex
%% Distinct proof detail retained from historical annex lines 121--130.

\noindent\textit{Reader note.}
This complete proof supports the separately authored range bridge in the body;
it does not replace any exact inherited proof.

\begin{annexitem}{annex:p04-universal.9}{thm:p04-numbering-partial-recursive}{Range of the local numbering}
\begin{proof}
If \(\varphi_e\simeq f\), Theorem~\ref{thm:enumeration}(a) makes \(f\)
partial computable, so the cited Theorem~\ref{thm:p04-tm-partial-recursive}
makes it partial recursive.  Conversely, that cited theorem makes every
partial recursive \(f\) partial computable, and
Theorem~\ref{thm:enumeration}(b) supplies a local index \(e\) with
\(\varphi_e\simeq f\).
\end{proof}
\end{annexitem}

\needspace{8\baselineskip}
\begin{annexitem}{annex:p04.local.counter-compilation}{thm:p04-counter-compilation}{Counter compilation}
\begin{proof}
We describe a uniform finite transition-table construction, including the
input/output interface.  Let the procedure use \(r\geq1\) counters.  On a
work interval of the tape store
\[
  \lhd\,\#1^{c_1}\#1^{c_2}\#\cdots\#1^{c_r}\#\,\rhd.
\]
The endpoint markers and separators are work symbols distinct from the
blank and numerical digits.  There are only finitely many additional
marked versions of these symbols.  The machine's control records the
current instruction, a scan mode, and, when needed, a number between
\(1\) and \(r\).  At an instruction boundary it scans the left marker.
The interval between the markers is finite, and every counter has exactly
its displayed tally block.  We now give the templates for each instruction.

To find counter \(i\), scan right from \(\lhd\), counting separators in
the finite control.  The next symbol is a separator exactly when the
counter is zero.  Otherwise scan to its closing separator; the preceding
symbol is its last tally.  For increment, mark that closing separator,
scan to \(\rhd\), and shift the suffix from the marked separator through
\(\rhd\) one cell to the right.  Carry out the shift from right to left:
remember the current symbol in the finite control, write it in its right
neighbor, and move to the next source cell on the left.  On reaching the
marked separator, write an ordinary separator to its right and a tally
in its old cell.  At each shift step the already processed suffix is in
its new position and the unprocessed prefix is unchanged.  Thus exactly
one tally has been inserted and all later counters are unchanged.

For a positive decrement, mark the last tally, and shift the suffix
beginning at the closing separator one cell to the left.  Process this
suffix from left to right, copying the next symbol into the hole and
moving the hole right until \(\rhd\) has been copied; blank its old cell.
The corresponding invariant is that the processed prefix has moved left
one cell and the remaining suffix is unchanged.  Exactly one tally is
deleted.  A zero test changes no cell.  Each completed instruction scans
back to \(\lhd\) and enters the control state for its selected target.
All these scans terminate because the marked interval is finite.

These are transition-table templates, not operations requiring an
unbounded control state.  In a scan mode there is one row for each member
of the fixed finite work alphabet; remembering a symbol introduces one
state for each possible symbol.  Separator counts range only up to the
fixed \(r\).  A stationary action is replaced by a move right to a fresh
state followed by a move left, with one preserving row for every symbol
in that fresh state.  Hence all moves have the required direction
\(L\) or \(R\).  Fresh states keep different instruction templates
disjoint.  Induction on executed instructions proves the boundary
invariant and exact agreement of counter values and control targets.

For completeness, initialize and return in the required binary convention.
Before using the work interval, put an origin marker on a spare track
at physical cell \(0\).  A track here means a finite product alphabet:
each tape cell stores an ordinary symbol together with one of finitely
many flags; a transition reads and writes that product symbol.  This
introduces finitely many symbols and still uses one head and one tape.
Copy the input digit string into a delimited work block and scan it from
left to right.  Starting at zero, for each digit \(b\) replace the input
counter \(c\) by \(2c+b\).  Doubling is obtained by moving its tallies
one at a time into a scratch counter, adding two tallies there for each
one removed, then moving the scratch counter back.  A one-bit flag in
the control decides the final increment.  These are finite sequences
of the instruction templates just described.  Copying the finite input
block can itself be performed by marking its next unprocessed digit,
carrying that digit in the control to the destination, appending it by
the suffix shift, and returning to the marked digit.  Preserve the origin
flag throughout shifts.  Initially the input is canonical, as required
by the numerical-machine interface.  The resulting counters are the
input value and zeros, and the procedure begins.

At a return, convert the selected tally block to binary.  Repeatedly
remove pairs of tallies, counting them in a scratch block; the possible
remaining tally is the remainder on division by two.  Record that bit
and repeat with the quotient until it is zero.  For an original zero
record the single digit zero.  Reverse the recorded bits by moving
their last digit successively into an output block.  This yields exactly
\(\operatorname{bin}(c_i)\).  Copy this block to the cells starting at
the preserved origin, write a blank immediately after it, remove work
flags from that prefix, and enter the accepting state.  To avoid an
overlap during this copy, first move the output block and workspace
past a blank interval longer than the output; finite scans determine
its length and perform the relocation.  No other tape cells or final
head location affect the output convention.  Every conversion terminates.
If the counter procedure never returns, its simulated instructions
continue forever, so the resulting Turing computation is infinite.

The procedure for printing this transition table is total.  Enumerate
the finite instruction list, instantiate the finitely many template
states and rows, and concatenate them with fresh state names.  The
templates for initialization and return are fixed; their counter indices
are finite syntactic data.  Literal initialization by \(a\) prints a
zeroing loop and \(a\) increment states, so printing terminates for every
natural number \(a\), irrespective of its cost.  Complete every unused
nonhalting state/symbol row by entering a fresh infinite-loop component.
Use distinct start, accepting, and rejecting states, even though rejection
is never used, and include the blank and both binary digits in the tape
alphabet and both digits in the input alphabet.

Finally assign numbers \(0,1,2\) to the designated states and symbols,
assign the remaining numbers by the order of their construction, sort
the full row list, and serialize its state count, alphabet size, input
bitmap and rows exactly as in
Definition~\ref{def:p04-program-serialization}.  Finite enumeration,
sorting and the specified tuple coding terminate and produce a valid
index.  Renaming invariance preserves all runs.  This proves both the
behavioral assertion and the total computability of the compiler.
\end{proof}
\end{annexitem}

\needspace{8\baselineskip}
\begin{annexitem}{annex:p04.local.universal}{thm:p04-local-universal}{Local interpreter}
\begin{proof}
Apply Theorem~\ref{thm:p04-counter-compilation} to the following counter
procedure.  Its numerical input is the
binary numeral for \(\operatorname{code}_2(e,x)\).  The following finite
arithmetic routines are available on counters.  A copy preserving a
counter moves its tallies to a target and a scratch counter in parallel,
then restores the source from the scratch counter.  Addition uses such
copies and repeated increment; multiplication uses repeated addition.
Comparison and truncated subtraction use copied counters, decrementing
them together until one is zero.  A loop with a copied bound executes
exactly that many iterations.  These routines terminate.  They implement
the bounded search for the inverse of the Cantor pairing already given
in Section~\ref{sec:turing-machines-01}; division by two is implemented
by removing pairs of tallies.  Consequently they also implement finite
binary digit extraction and right-nested tuple encoding and decoding.
Every routine has a fixed finite instruction list, using scratch counters
which it clears before reuse.

There is no bound on the length of a simulated program or tape.  Variable
finite lists are stored as two numbers: their length and their nested
pair code.  The empty list has length zero and code zero; a nonempty list
with first entry \(a\) and remaining list code \(v\) uses
\(\operatorname{pair}_{\N}(a,v)\), with a terminal zero after its last
entry.  This auxiliary list code is separate from the source's
\(\operatorname{code}_k\), whose last entry is unpaired.  Length-bounded
pairing inversions implement scan, lookup, and reconstruction of these
lists.  Thus unbounded storage size uses counter values, not additional
counters or control states.

Decode \((e,x)\), invert the outer pairing of \(e\), and apply the exact
finite validity checks in
Definition~\ref{def:p04-program-serialization}: the stated length,
\(n,b\geq3\), input bitmap, all row-field bounds, directions and ordered
nonhalting keys.  Reject an incorrect length before interpreting its
fields.  Each check uses finite arithmetic and a length-bounded scan.
If it fails, enter the one-instruction infinite increment loop.  Otherwise
retain the full decoded transition table as a list of five-field rows.

Represent a simulated configuration by its state \(q\), its head
coordinate \(h\), and a finite dictionary of its nonblank cells.  Encode
integers by \(2j\) for \(j\geq0\) and \(2(-j)-1\) for \(j<0\);
left and right movement are terminating arithmetic operations on these
codes.  A dictionary is a list of coordinate/symbol pairs with distinct
coordinates; an absent coordinate means blank.  Lookup scans its finite
length.  Replacement changes its matching entry, deletes it when the
new symbol is blank, or appends a previously absent nonblank entry.
These operations preserve distinctness.  Initialize \(q=0,h=0\) and
the cells \(0,\ldots,|\operatorname{bin}(x)|-1\) with the corresponding
designated digit numbers \(1,2\).

At a nonhalting state, look up the symbol at \(h\), find its unique row
in the full transition table, replace that cell, move the head one cell
in the row's direction, and update the state.  Each such operation
terminates, because the table and current dictionary are finite.  By
induction on the number of simulated steps, the state, head and dictionary
represent exactly the configuration of the decoded machine after that
many steps.  The base case is the specified binary input; the induction
step is precisely the defining one-step rule.  Thus an infinite machine
run gives infinitely many completed simulation steps, and every finite
halt is reached after finitely many interpreter steps.

At state \(2\), enter the infinite loop.  At state \(1\), scan coordinates
\(0,1,\ldots\) until the first blank; this scan terminates since the
dictionary is finite.  Check that the prefix is nonempty, uses only the
designated binary digits, and has no leading zero unless it consists of
the single zero digit.  An invalid prefix leads to the infinite loop.
For a valid prefix evaluate its binary value by the recurrence
\(v\mapsto2v+b\), and return it.  The compiler then writes its canonical
binary output.  This is exactly the definition of \(\varphi_e(x)\),
including every undefined case.  The compiled procedure therefore
computes the asserted partial function \(U\).
\end{proof}
\end{annexitem}

\needspace{8\baselineskip}
\begin{annexitem}{annex:p04.local.parameter}{thm:p04-local-parameter}{Literal specialization}
\begin{proof}
Fix \(m,n\geq1\).  Given numbers \(e,a_1,\ldots,a_m\), print a counter
procedure with those literal initializations and the following fixed
instruction blocks.  On numerical input \(z\), decode
\(z=\operatorname{code}_n(x_1,\ldots,x_n)\), compute
\[
  w=\operatorname{code}_{m+n}(a_1,\ldots,a_m,x_1,\ldots,x_n),
\]
and execute the counter interpreter in
Annex~\ref{annex:p04.local.universal} on the pair \((e,w)\).
For \(n=1\) decoding is the identity; for larger \(n\) it uses exactly
\(n-1\) pairing inversions.  The two blocks for coding and decoding are
finite terminating counter procedures by that annex.  Supply the
interpreter's input counters and enter its start instruction.  Rename all
scratch counters and instruction labels beforehand so that literal
storage and runtime data are not overwritten by a different block.
Clear each scratch counter at its first use.  This is a mechanical
operation on finite instruction lists.

At the three interfaces, respectively, the designated counters contain
\((e,\bar a,z)\), then \((e,\bar a,\bar x)\), and then \((e,w)\).
Termination of the first two blocks proves these assertions for every
input; the interpreter invariant proves that the last block returns
exactly when \(\varphi_e(w)\) is defined, with that value.  The
\((m+n)\)-ary application convention identifies this value with
\(\varphi_e(\bar a,\bar x)\).  In particular, if \(e\) is invalid the
generated program reaches the interpreter's infinite loop rather than
failing to compile.

Apply the total compiler from
Theorem~\ref{thm:p04-counter-compilation} to this printed procedure and
call its valid index \(\mathsf{par}_n^m(e,\bar a)\).  Production of
literal chains, finite instruction concatenation, fresh names,
normalization and serialization all terminate.  Thus the index map is
total computable, independently of whether its output program returns.
On input \(\operatorname{bin}(\operatorname{code}_n(\bar x))\) the
compiled machine has exactly the behavior just proved, establishing
the required strong equality.  If \(m,n\) are also supplied, a loop
prints their finitely many tuple and initialization blocks; the same
construction is total on positive arities.  No search for an index of
a semantically specified function occurs in this compiler.
\end{proof}
\end{annexitem}

\needspace{8\baselineskip}
\begin{annexitem}{annex:p04.local.recursion}{thm:p04-local-recursion}{Diagonal specialization}
\begin{proof}
Fix \(k\geq1\) and the partial computable function \(g\).  Choose a
valid index \(b\) of arity \(k+1\) for \(g\) by
Theorem~\ref{thm:enumeration}.  Put
\[
  d(z)=\mathsf{par}_k^1(z,z),
  \qquad
  H(z,\bar x)\simeq g(d(z),\bar x).
\]
The map \(d\) is total computable by
Theorem~\ref{thm:p04-local-parameter}.  More explicitly, compute
\(d(z)\), form the tuple code of \((d(z),\bar x)\), and run
\(U(b,\operatorname{code}_{k+1}(d(z),\bar x))\).
The literal compiler and the interpreter give a machine for this
composition, so \(H\) is partial computable even when \(g\) is not
total.  Choose a valid index \(a\) for \(H\) of arity \(k+1\), and
let \(e=d(a)\).  This is a valid index because the parameter compiler
always produces one.  For every \(\bar x\),
\[
 \varphi_e(\bar x)
 \simeq\varphi_a(a,\bar x)
 \simeq H(a,\bar x)
 \simeq g(d(a),\bar x)
 =g(e,\bar x).
\]
The first equality is parameter specialization, the second is the chosen
index property, and the third is the definition of \(H\).  Each
\(\simeq\) equates both definedness and value.  This is Kleene's
diagonal argument with the local specialization map in place of his
equation-code operation.

For a total computable \(f\), take \(k=1\) and
\(g(z,x)\simeq U(f(z),x)\).  Its machine first computes \(f(z)\),
which always returns, and then runs the interpreter.  The proved result
gives \(\varphi_e(x)\simeq U(f(e),x)\simeq\varphi_{f(e)}(x)\).
An invalid value \(f(e)\) simply denotes the everywhere undefined
function, as required by the fixed numbering.

The particular intermediate transformation used in the inherited sketch
can now also be constructed.  The procedure
\(B(x,y)\simeq U(U(x,x),y)\) first runs the inner interpreter and
starts the outer one only if the inner call returns.  Let \(c\) be a
valid binary index for this procedure, and set
\(p(x)=\mathsf{par}_1^1(c,x)\).  This map is total computable and
satisfies
\(\varphi_{p(x)}(y)\simeq\varphi_{\varphi_x(x)}(y)\), with both
sides undefined when the inner computation fails to return.  Thus the
inherited fixed-point calculation has its previously omitted program
transformation, without invoking an additional recursion theorem.
\end{proof}
\end{annexitem}
